% Experimental Methodology chapter for:
% A Hierarchical Time-Aware Transformer for Flow-Level Network Intrusion Detection
%
% Required packages already used by the thesis:
%   amsmath, booktabs, multirow, array, graphicx
% Recommended in the main preamble:
%   \usepackage{amssymb}
%
% The platform and purchased allocation are known. Replace the remaining
% environment macros with values captured from each final Colab runtime.
% In particular, do not replace \ExpGPU with only "GPU": record T4 or L4.
\providecommand{\ExpPlatform}{Google Colaboratory Pro+}
\providecommand{\ExpComputeBudget}{600 purchased compute units}
\providecommand{\ExpGPU}{NVIDIA T4 (16~GB) or NVIDIA L4 (24~GB); exact model recorded per run}
\providecommand{\ExpCPU}{\textit{[insert CPU model and RAM]}}
\providecommand{\ExpOS}{Google Colab hosted Linux runtime; \textit{[insert runtime image details]}}
\providecommand{\ExpSoftware}{\textit{[insert Python, PyTorch, CUDA, and scikit-learn versions]}}
\providecommand{\ExpSuricata}{\textit{[insert Suricata version and custom-plugin revision]}}

\chapter{Experimental Methodology}
\label{sec:experimental_methodology}

This chapter defines the protocol used to evaluate the proposed hierarchical
detector. Its purpose is to separate model development from final evaluation,
to make comparisons between alternatives fair, and to ensure that every
reported number can be traced to a dataset manifest, model configuration,
checkpoint, decision threshold, and random seed. This separation is especially
important in intrusion detection because temporal leakage, class imbalance,
and repeated use of the test set can produce highly optimistic conclusions
\cite{SommerPaxson2010,Arp2022DosDonts}.

To avoid duplicating earlier chapters, this chapter does not redefine the raw
packet and flow schemas, the mathematical detection problem, or the internal
layers of the proposed models. Dataset provenance and transformations are
given in Section~\ref{sec:dataset_preprocessing}, the research questions and
notation in Section~\ref{sec:problem_formulation}, and the architecture in
Section~\ref{sec:proposed_methodology}. The present chapter specifies only the
experimental comparisons, controls, optimization protocol, metrics,
statistical analysis, and computational measurement needed to test that
methodology.

The evaluation is organized around the two research questions mapped to
controlled comparisons in Table~\ref{tab:rq_mapping}. Stage~1 experiments
address whether ordered packet records,
continuous-time information, and flow statistics improve the representation
of an individual flow. Stage~2 experiments address whether causal context from
earlier related flows improves the classification of the current flow. In both
stages, macro-F1 is the primary threshold-dependent metric. Macro-precision,
macro-recall, class-1 performance, ROC-AUC, PR-AUC, the confusion matrix, and
the false-positive rate (FPR) provide complementary evidence. Parameter count,
training time, and inference time are reported to distinguish predictive gains
from increases in computational cost.

\section{Experimental Logic}
\label{subsec:experimental_questions}

The research-question definitions and their conceptual controls already appear
in Table~\ref{tab:rq_mapping}; they are not repeated here. They are converted
below into named experiment groups. Within a group, one factor is changed while
the ordered data manifest, fitted preprocessor, optimization budget, model
selection rule, and evaluation code remain fixed. Stage~1 groups test the
intra-flow representation, Stage~2 groups test the additional value of causal
inter-flow history, and a final measurement layer tests statistical stability
and computational cost. This design prevents a gain caused by a different
split, threshold, or training rule from being attributed to the architecture.

Experiments are divided into \emph{exploratory} and \emph{confirmatory}
comparisons. Hyperparameter and sensitivity sweeps are exploratory and use the
validation partition to identify useful settings. Confirmatory comparisons use
the resulting frozen configuration and evaluate the test partition once. The
test partition is never used to choose sequence length, context method, window
size, checkpoint, loss, or threshold. This distinction follows established
recommendations for avoiding test-set overfitting in empirical machine
learning and security research \cite{Arp2022DosDonts,Pineau2021Reproducibility}.

\section{Dataset Partition and Leakage Controls}
\label{subsec:experimental_data_protocol}

All experiments use the aligned flow manifest constructed in
Section~\ref{sec:dataset_preprocessing}. The PCAP provenance, exact packet and
flow counts, label distribution, selected predictors, and exclusions are
reported once in Sections~\ref{subsec:raw_composition}--
\ref{subsec:feature_selection} and are not restated here. The experimental
invariant is that every comparison uses the same flow universe and the same
training-fitted preprocessing state unless a table explicitly identifies a
separate exploratory manifest.

\subsubsection{Chronological Hold-Out Protocol}
\label{subsubsec:chronological_holdout}

The fixed 70/10/20 chronological split and its exact counts are defined in
Table~\ref{tab:chronological_split}. Unique flows are stably ordered by
\texttt{flow\_start\_timestamp\_us}, with \texttt{flow\_id} as the
deterministic tie-breaker, and all packets from one flow remain in the same
partition. This chronological design evaluates the more deployment-like
setting in which a detector is trained on earlier traffic and applied to later
traffic, rather than allowing temporally adjacent records to be distributed
randomly across partitions \cite{SommerPaxson2010,Arp2022DosDonts}.

All data-dependent transformations are fitted using the training partition
only. These include categorical vocabularies, quantile clipping bounds,
logarithmic transformations, and robust-scaling parameters. The fitted
preprocessor is then applied without refitting to validation and test data.
Class weights and sampling probabilities are also calculated from training
labels only.

\subsubsection{Stage~2 Causality and Split Policy}
\label{subsubsec:stage2_evaluation_causality}

Stage~2 consumes embeddings exported by one frozen Stage~1 checkpoint. Its
context builder scans flows in chronological order and updates a history only
after the current example has been constructed. Consequently, the context of
flow $i$ may contain only flows with earlier timestamps; neither future
embeddings nor historical labels are retrieved. Under the principal
\texttt{online} policy, a validation example may use earlier training flows,
and a test example may use earlier training, validation, and test flows. This
simulates a continuously operating detector whose past observations remain
available, while preserving the rule that no future observation can influence
the present prediction.

All Stage~2 model comparisons use exactly the same Stage~1 embedding files,
metadata ordering, context indices, and left-padding masks. The Stage~1 encoder
is not fine-tuned during Stage~2 training. A change in Stage~1 checkpoint would
therefore constitute a new experiment rather than a Stage~2 architecture
ablation.

\paragraph{Comparability rule.}
Only runs evaluated on the same ordered flow manifest may appear in the same
ranking table or paired statistical test. Several early Stage~1 experiments
produced results for $35{,}410$ test flows, whereas the principal full-data
manifest contains $84{,}380$ test flows. These results measure performance on
different empirical distributions and are not directly comparable. The final
main tables must therefore either rerun all relevant Stage~1 alternatives on
the $84{,}380$-flow manifest or present the smaller-manifest experiments in a
separate exploratory table with their sample count stated explicitly.

\section{Experimental Matrix}
\label{subsec:experimental_matrix}

The complete evaluation matrix is summarized in
Table~\ref{tab:experiment_matrix}. The matrix separates comparisons intended
to evaluate the packet encoder from comparisons intended to evaluate the full
hierarchical system.

\begin{table}[t]
    \centering
    \caption{Experimental matrix. Only the factor listed in the final column is
    changed within each experiment group.}
    \label{tab:experiment_matrix}
    \resizebox{\textwidth}{!}{%
    \begin{tabular}{@{}lllll@{}}
        \toprule
        Group & Stage & Fixed reference & Compared alternatives & Varied factor \\
        \midrule
        S1-E1 & 1 & Full chronological split & Flow MLP; packet
        sequence models; proposed Stage~1 & Encoder family \\
        S1-E2 & 1 & Scheme B, $L=128$, head & Both, position-only,
        time-only, no encoding & Temporal/order encoding \\
        S1-E3 & 1 & Both encodings, $L=128$, head & Schemes A, B, C &
        Flow-statistics integration \\
        S1-E4 & 1 & Scheme B, head & $L=8,16,32,64,128,256$ & Packet budget \\
        S1-E5 & 1 & Scheme B, $L=128$ & Head, head-tail, random &
        Truncation policy \\
        S2-E1 & 2 & Same frozen Stage~1 embeddings & No-context MLP;
        Transformer; LSTM; GRU; CNN+LSTM; proposed & Context encoder \\
        S2-E2 & 2 & Proposed model, $W=128$ & Time-only, source-host,
        destination-host, endpoint & Context relation \\
        S2-E3 & 2 & Proposed model, source-host & $W=16,32,64,128,256$ &
        Context capacity \\
        \bottomrule
    \end{tabular}%
    }
\end{table}

\subsubsection{Stage~1 Baselines}
\label{subsubsec:stage1_baselines}

The Stage~1 comparison includes the following baselines:

\begin{itemize}[leftmargin=*]
    \item \textbf{Flow-level MLP.} A multilayer perceptron consumes only the
    aggregated statistical flow vector. It provides the direct reference for
    testing whether packet-level sequence information adds value beyond
    conventional flow aggregation.

    \item \textbf{CNN1D.} Parallel one-dimensional convolutions with kernel
    sizes 3, 5, and 7 capture local packet motifs, followed by global pooling
    and classification. A second variant concatenates the log inter-arrival
    feature to each packet record. The design follows the general use of
    multi-kernel convolution for sequence classification \cite{Kim2014CNN}.

    \item \textbf{LSTM and BiLSTM.} Two-layer recurrent encoders with hidden
    width 128 and mean pooling provide unidirectional and bidirectional
    sequence references. The LSTM-with-time variant concatenates the stored
    log inter-arrival value to each packet vector
    \cite{Hochreiter1997LSTM}.

    \item \textbf{GRU.} A two-layer, width-128 GRU with mean pooling tests
    whether a lighter gated recurrent encoder is sufficient
    \cite{Cho2014GRU}.

    \item \textbf{Standard Transformer.} This control has the same packet
    projection and self-attention family as the proposed Stage~1 model but
    uses conventional order encoding without the proposed cumulative-time
    mechanism \cite{Vaswani2017}.
\end{itemize}

The recurrent and convolutional baselines use the same packet feature tensors,
sequence masks, chronological split, focal objective, sampler, optimizer
budget, validation-based checkpoint rule, and threshold-selection rule as the
proposed model \emph{within the Stage~1 model-family benchmark}. In the current
comparison trainer, every member of this benchmark, including the proposed
encoding variants, is checkpointed by validation class-1 F1 and obtains its
threshold from the validation precision-recall points. This shared rule makes
the rows within that benchmark comparable. It differs from the standalone
principal Stage~1 trainer described below, which checkpoints by validation
PR-AUC. The two protocols are therefore kept in separate tables unless all
models are rerun through one common trainer. Model capacity is not forced to be
identical because doing so can make one architecture artificially weak;
instead, trainable parameter count and runtime are reported next to predictive
performance.

\subsubsection{Stage~1 Ablations and Sensitivity Tests}
\label{subsubsec:stage1_ablation_protocol}

Four encoding variants isolate the contribution of order and elapsed time:
(i) both positional and time-aware encoding, (ii) positional encoding only,
(iii) time-aware encoding only, and (iv) neither encoding. These variants keep
the packet features and Transformer capacity unchanged. Their comparison is
particularly important because concatenating an interval value to a packet
vector is not equivalent to using elapsed time as the coordinate of a
sinusoidal sequence encoding.

Flow-statistics integration is evaluated through the three schemes defined in
Table~\ref{tab:stage1_fusion_variants}. Scheme A repeats flow statistics at
every packet position, Scheme B excludes flow statistics, and Scheme C uses a
separate flow branch with bounded token-FiLM conditioning. Scheme B is used
for packet-encoder and truncation studies because it isolates sequence
modelling from flow-statistics fusion. Scheme C is the principal complete
Stage~1 model and supplies the embeddings used by Stage~2.

Sequence-length sensitivity is measured with
$L\in\{8,16,32,64,128,256\}$ under \texttt{head} selection. Selection-policy
sensitivity fixes $L=128$ and compares \texttt{head},
\texttt{head\_tail}, and deterministic \texttt{random} subsampling. All other
settings remain unchanged. In particular, a new threshold is selected on the
validation predictions of each trained variant; the test set is not searched
for a favorable operating point.

\subsubsection{Stage~2 Baselines and Context Ablations}
\label{subsubsec:stage2_baseline_protocol}

The proposed target-query gated attention model is compared with five controls:

\begin{itemize}[leftmargin=*]
    \item \textbf{No-context MLP} classifies only the current Stage~1
    embedding and therefore measures the gain attributable to historical
    context.

    \item \textbf{Vanilla Transformer} applies full self-attention to the same
    context sequence and classifies its last valid token
    \cite{Vaswani2017}.

    \item \textbf{LSTM} and \textbf{GRU} replace attention with recurrent
    context encoders while preserving the Stage~1 inputs, context window,
    masks, and classifier protocol \cite{Hochreiter1997LSTM,Cho2014GRU}.

    \item \textbf{CNN+LSTM} applies parallel temporal convolutions with kernel
    sizes 3, 5, and 7 before a two-layer LSTM, combining local pattern
    extraction with recurrent context modelling
    \cite{Kim2014CNN,Hochreiter1997LSTM}.
\end{itemize}

For the architecture comparison, contextual baselines receive the same
\texttt{source\_host} histories with $W=128$ and the current flow included as
the last valid element. With the target included, the maximum historical
capacity is $W-1=127$ flows. The no-context model receives the same target
embedding but discards all preceding elements.

Context-relation sensitivity fixes the proposed encoder and $W=128$, then
changes only the eligibility relation among \texttt{time\_only},
\texttt{source\_host}, \texttt{destination\_host}, and \texttt{endpoint}.
Window-size sensitivity fixes \texttt{source\_host} and evaluates
$W\in\{16,32,64,128,256\}$. For each context configuration, the dataset report
records the minimum, median, mean, 95th percentile, and maximum number of valid
elements. This reveals whether a nominally larger window actually supplies
more history or is mostly padding for low-activity hosts.

\section{Training, Model Selection, and Thresholding}
\label{subsec:experimental_training_protocol}

\subsubsection{Principal Training Configuration}
\label{subsubsec:principal_training_config}

The architecture and optimization settings follow
Table~\ref{tab:principal_method_config}. The main values relevant to
experimental replication are summarized in
Table~\ref{tab:experimental_training_config}. AdamW is used in both stages to
decouple weight decay from the adaptive gradient update
\cite{Loshchilov2019AdamW}. Focal loss reduces the influence of easy examples
from the majority class \cite{Lin2017FocalLoss}.

\begin{table}[t]
    \centering
    \caption{Principal optimization configuration. Values changed in a named
    ablation are the only exceptions.}
    \label{tab:experimental_training_config}
    \begin{tabular}{@{}p{3.0cm}p{4.6cm}p{5.7cm}@{}}
        \toprule
        Setting & Stage~1 & Stage~2 \\
        \midrule
        Architecture and input & Table~\ref{tab:principal_method_config} &
        Table~\ref{tab:principal_method_config} \\
        Batch size & 128 & 128 \\
        Optimizer & AdamW & AdamW \\
        Learning rate & $1\times10^{-4}$ & $3\times10^{-5}$ \\
        Weight decay & $1\times10^{-4}$ & $3\times10^{-4}$ \\
        Gradient clipping & Global norm 1.0 & Global norm 0.5 \\
        Loss & Focal, $\gamma=1$ & Focal, $\gamma=1$ \\
        Sampler target & 9\% class-1 examples & 5\% class-1 examples \\
        Maximum epochs & 500 & 500 \\
        Early-stopping patience & 25 epochs & 15 epochs \\
        Checkpoint criterion & Validation PR-AUC & Validation PR-AUC \\
        Random seed & 130 & 130 \\
        \bottomrule
    \end{tabular}
\end{table}

The weighted sampler is used only to construct training mini-batches and draws
with replacement. Validation, test, and embedding-export loaders are
non-sampled and preserve one row per flow. This distinction prevents the
sampler from duplicating flow embeddings or altering the empirical label
distribution used for evaluation. Automatic mixed precision is enabled on
CUDA. Stage~1 reduces the learning rate when validation PR-AUC plateaus;
Stage~2 uses five warm-up epochs followed by cosine annealing.

\subsubsection{Validation-Only Model Selection}
\label{subsubsec:validation_only_selection}

For every configuration, early stopping and checkpoint selection use only the
validation partition. The standalone principal Stage~1 pipeline and the
Stage~2 pipeline restore the checkpoint with the highest validation PR-AUC.
PR-AUC is appropriate for these principal runs because it evaluates ranking
quality across decision thresholds and is more informative than accuracy under
strong class imbalance \cite{SaitoRehmsmeier2015}. As noted in
Section~\ref{subsubsec:stage1_baselines}, the current Stage~1 model-family
benchmark instead uses validation class-1 F1 for all rows. This difference is
reported rather than hidden: comparisons are made within a common protocol,
and a row from the PR-AUC-selected standalone run is not inserted into the
class-1-F1-selected benchmark table. A unified confirmatory rerun should use
validation PR-AUC for every model.

Hyperparameters are selected in two phases. First, pilot experiments determine
reasonable optimization ranges using training and validation data. Second,
the chosen configuration is frozen before the confirmatory test run. When an
ablation changes a structural factor, it receives the same optimization budget
as the reference model; it is not tuned on the test result. All attempted
configurations, including unsuccessful ones, should remain in the run manifest
to make the extent of validation-set search transparent
\cite{Pineau2021Reproducibility,Bouthillier2021Variance}.

\subsubsection{Decision-Threshold Selection}
\label{subsubsec:experimental_thresholding}

Let $p_i$ denote the restored model's probability for class 1. A prediction at
threshold $\tau$ is
\begin{equation}
    \widehat{y}_i(\tau)=\mathbb{I}[p_i\geq\tau].
    \label{eq:experimental_decision_rule}
\end{equation}
The threshold is selected by maximizing validation class-1 F1:
\begin{equation}
    \tau^{*}=\arg\max_{\tau\in\mathcal{T}}
    F_{1,1}^{\mathrm{val}}(\tau).
    \label{eq:validation_threshold_selection}
\end{equation}
The standalone Stage~1 trainer searches
$\mathcal{T}=\{0.100,0.101,\ldots,0.949\}$, whereas Stage~2 uses 199 evenly
spaced values from 0.01 to 0.99. The Stage~1 model-family comparison trainer
evaluates the exact score thresholds returned by the validation
precision-recall curve. Although the candidate sets differ between experiment
blocks, every selected threshold is based on validation predictions only and
is frozen before test evaluation. ROC-AUC and PR-AUC do not depend on this
threshold. A threshold optimized on test labels may be retained only as a
clearly marked oracle diagnostic for measuring threshold-transfer gap; it is
not an official result and must not be used in model comparisons.

\section{Evaluation Metrics}
\label{subsec:experimental_metrics}

The implementation uses the metric definitions provided by scikit-learn
\cite{Pedregosa2011ScikitLearn}. Label 1 is the alert-derived positive class and
label 0 is the negative class. The confusion matrix is always reported in the
following orientation:
\begin{equation}
    \mathbf{C}=
    \begin{bmatrix}
        \mathrm{TN} & \mathrm{FP}\\
        \mathrm{FN} & \mathrm{TP}
    \end{bmatrix},
    \label{eq:experimental_confusion_matrix}
\end{equation}
where rows denote true labels and columns denote predicted labels. Stating the
orientation prevents an ambiguity that is common when confusion matrices are
copied between software packages.

\subsubsection{Class-Wise and Macro-Averaged Measures}
\label{subsubsec:macro_metrics}

For class $c\in\{0,1\}$, precision, recall, and F1 are
\begin{align}
    P_c &= \frac{\mathrm{TP}_c}
    {\mathrm{TP}_c+\mathrm{FP}_c}, \\
    R_c &= \frac{\mathrm{TP}_c}
    {\mathrm{TP}_c+\mathrm{FN}_c}, \\
    F_{1,c} &= \frac{2P_cR_c}{P_c+R_c}.
    \label{eq:classwise_metrics}
\end{align}
Undefined ratios are assigned zero, matching the implementation's
\texttt{zero\_division=0} setting. The macro-averaged metrics are
\begin{align}
    P_{\mathrm{macro}} &= \frac{P_0+P_1}{2}, \\
    R_{\mathrm{macro}} &= \frac{R_0+R_1}{2}, \\
    F_{1,\mathrm{macro}} &= \frac{F_{1,0}+F_{1,1}}{2}.
    \label{eq:experimental_macro_metrics}
\end{align}

Macro-F1 is the primary F1 reported in this thesis and is denoted simply as
``F1'' in high-level result tables. It weights benign and positive classes
equally despite their different frequencies. Macro-precision and macro-recall
are the corresponding primary precision and recall measures. Because macro
averaging can conceal which error type changed, $P_1$, $R_1$, and $F_{1,1}$
are always reported alongside the macro values.

Accuracy and weighted averages are included for completeness but are treated
as secondary. On the raw dataset, an always-negative classifier would achieve
approximately $96.055\%$ accuracy while obtaining zero positive recall. High
accuracy or weighted F1 is therefore not sufficient evidence of an effective
intrusion detector.

\subsubsection{False-Positive Rate and Operational Error Counts}
\label{subsubsec:fpr_metric}

The false-positive rate is
\begin{equation}
    \mathrm{FPR}=\frac{\mathrm{FP}}{\mathrm{FP}+\mathrm{TN}},
    \label{eq:false_positive_rate}
\end{equation}
and the false-negative rate is
\begin{equation}
    \mathrm{FNR}=\frac{\mathrm{FN}}{\mathrm{FN}+\mathrm{TP}}
    =1-R_1.
    \label{eq:false_negative_rate}
\end{equation}
FPR is emphasized because even a small rate can create a large alert volume in
a high-throughput network. Each result table therefore reports both FPR and the
absolute numbers of FP and FN. For a test partition containing $N_0$ benign
flows, the expected number of false alerts at the evaluated operating point is
$N_0\times\mathrm{FPR}$.

\subsubsection{Threshold-Independent Ranking Measures}
\label{subsubsec:ranking_metrics}

ROC-AUC summarizes the true-positive rate as the false-positive rate varies
over all score thresholds \cite{Fawcett2006ROC}. It measures ranking quality
but can remain high when the negative class dominates. PR-AUC focuses on the
positive class and is therefore given greater weight for this imbalanced
dataset \cite{SaitoRehmsmeier2015}.

The code field named \texttt{pr\_auc} is computed with scikit-learn's
\texttt{average\_precision\_score}. It is therefore average precision (AP),
not trapezoidal integration of linearly interpolated precision-recall points:
\begin{equation}
    \mathrm{AP}=\sum_{k}(R_k-R_{k-1})P_k,
    \label{eq:average_precision}
\end{equation}
where $P_k$ and $R_k$ are the precision and recall at successive score
thresholds. To remain consistent with the stored result files, this thesis
uses the label ``PR-AUC (AP)'' on first mention and ``PR-AUC'' thereafter. A
random ranking has expected AP close to the class-1 prevalence, so the observed
value should be interpreted relative to the test partition's positive rate.

\section{Statistical Reliability}
\label{subsec:statistical_reliability}

Neural results can vary with initialization and mini-batch order; a single
best run can therefore overstate an architecture's advantage
\cite{ReimersGurevych2017,Bouthillier2021Variance}. Existing exploratory
sequence-length, packet-selection, context-relation, and context-window sweeps
use seed 130 and are explicitly labelled as single-seed sensitivity analyses.
The confirmatory Stage~1 and Stage~2 model comparisons are repeated with five
seeds, $\{130,131,132,133,134\}$. Each seed uses the same immutable data split
and hyperparameters. The results chapter reports mean and standard deviation
for macro-F1, PR-AUC, ROC-AUC, FPR, and class-1 F1. This distinction preserves
the useful completed sweeps without presenting one stochastic run as an
estimate of training stability.

Two complementary paired analyses are applied to the fixed test predictions.
First, a stratified paired bootstrap draws, with replacement, the original
number of class-0 and class-1 test indices. The same sampled indices are used
for both models in every replicate. For metric $M$, replicate $b$ yields
\begin{equation}
    \Delta_b=M(\mathcal{S}_b,A)-M(\mathcal{S}_b,B).
    \label{eq:paired_bootstrap_difference}
\end{equation}
With $B=1{,}000$ replicates, the 2.5th and 97.5th percentiles of
$\{\Delta_b\}_{b=1}^{B}$ form a 95\% paired percentile interval
\cite{EfronTibshirani1993Bootstrap}. This analysis is applied to macro-F1,
PR-AUC, ROC-AUC, and FPR. It quantifies uncertainty conditional on the observed
test flows; it does not measure cross-network generalization or replace the
multi-seed analysis.

Second, hard classification errors are compared with McNemar's test because
the two models predict the same flows \cite{McNemar1947}. If $n_{01}$ is the
number of flows classified correctly only by model A and $n_{10}$ is the
number classified correctly only by model B, the continuity-corrected statistic
is
\begin{equation}
    \chi^2=\frac{\left(|n_{01}-n_{10}|-1\right)^2}
    {n_{01}+n_{10}}.
    \label{eq:mcnemar_statistic}
\end{equation}
The significance level is $\alpha=0.05$. When several alternatives are tested
against the proposed model, Holm's step-down procedure adjusts the resulting
$p$-values to control family-wise error \cite{Holm1979}. Statistical
significance is reported together with the observed metric difference and its
confidence interval; a small $p$-value alone is not treated as evidence of a
practically important gain.

\section{Efficiency and Resource Measurement}
\label{subsec:efficiency_protocol}

Experiments are executed in Google Colaboratory Pro+ using a company-funded
allocation of 600 compute units. Colab compute units represent an access and
usage budget, not a fixed amount of GPU work, and Colab states that accelerator
availability and runtime limits may change over time
\cite{GoogleColabFAQ2026}. The allocation is therefore documented as a resource
budget but is not used as an efficiency metric.

The available accelerators are an NVIDIA T4 with 16~GB of memory or an NVIDIA
L4 with 24~GB \cite{NVIDIAT4Datasheet,NVIDIAL4ProductBrief}. The exact GPU
assigned to every run is captured from the live runtime. For each
confirmatory seed, all models in a comparison block are run on the same GPU
family. Exploratory predictive results produced on T4 and L4 may be retained
when their hardware is disclosed, but paired significance claims require a
hardware-matched rerun. Wall-clock times from T4 and L4 are never directly
ranked. Every timed comparison block uses one GPU type with the same software
environment, batch size, input manifest, and precision mode. If Colab
reassigns a different accelerator, that run starts a separate timing block.
Runs interrupted by runtime termination or resumed in a new virtual machine
remain usable for predictive evaluation when the checkpoint is intact, but
they are excluded from wall-clock efficiency comparisons unless all timed
segments are explicitly recorded and summed.
This constraint is necessary because accelerator type, CUDA kernels, library
versions, and deterministic settings materially affect runtime
\cite{Pineau2021Reproducibility}.

\begin{table}[t]
    \centering
    \caption{Execution environment for the final experiments. Dynamic values
    are copied from the runtime that produced each reported timing block.}
    \label{tab:experimental_environment}
    \begin{tabular}{@{}p{3.4cm}p{10.0cm}@{}}
        \toprule
        Component & Recorded environment \\
        \midrule
        Hosted platform & \ExpPlatform \\
        Purchased allocation & \ExpComputeBudget \\
        Accelerator & \ExpGPU \\
        Host processor and RAM & \ExpCPU \\
        Operating environment & \ExpOS \\
        Learning software & \ExpSoftware \\
        Traffic extractor & \ExpSuricata \\
        \bottomrule
    \end{tabular}
\end{table}

\subsubsection{Parameter Count and Model Size}
\label{subsubsec:parameter_measurement}

For a PyTorch model with parameter tensors $\{\theta_j\}$, total and trainable
parameter counts are
\begin{align}
    N_{\mathrm{total}} &= \sum_j |\theta_j|, \\
    N_{\mathrm{train}} &= \sum_{j:\,\theta_j\ \mathrm{trainable}}|\theta_j|.
    \label{eq:parameter_counts}
\end{align}
The principal complexity column reports $N_{\mathrm{train}}$; total and
non-trainable counts are retained in the run artifact. The approximate FP32
parameter footprint is
\begin{equation}
    S_{\mathrm{FP32}}=\frac{4N_{\mathrm{total}}}{1024^2}\ \mathrm{MiB},
    \label{eq:fp32_model_size}
\end{equation}
excluding optimizer states, activations, and input batches. Parameter count is
reported separately for Stage~1 and Stage~2. The end-to-end hierarchy may also
report their sum, but Stage~1 and Stage~2 training times must not be added to
claim online latency because training is an offline operation.

\subsubsection{Training Time}
\label{subsubsec:training_time_measurement}

Training time is measured with a monotonic wall-clock timer around the complete
fit loop. It includes forward and backward passes, validation at the end of
each epoch, validation-threshold search when enabled, checkpoint writing, and
early stopping. It excludes PCAP extraction, CSV preprocessing, precomputed
tensor construction, Stage~1 embedding export, and Stage~2 context-index
construction. The number of completed epochs and average seconds per completed
epoch are reported with total time, because early stopping can make two models
with identical epoch budgets appear to have different computational cost.

Stage~1 and Stage~2 times are compared only within their respective stage and
accelerator block. The
two fit loops perform different validation and artifact-writing work, so their
absolute times do not represent a controlled cross-stage benchmark. On CUDA,
device synchronization is performed at timing boundaries wherever asynchronous
kernel execution could otherwise cause under-measurement.

\subsubsection{Inference Time, Throughput, and Scope}
\label{subsubsec:inference_time_measurement}

Inference is executed with gradients disabled and the model in evaluation
mode. The reported test inference time, $T_{\mathrm{test}}$, covers one full
non-shuffled pass through the test DataLoader, including host-to-device
transfer, forward computation, probability collection, and metric computation.
CUDA is synchronized immediately before and after the timed call. From
$N_{\mathrm{test}}$ evaluated flows, per-flow time and throughput are
\begin{align}
    t_{\mathrm{flow}} &= \frac{T_{\mathrm{test}}}{N_{\mathrm{test}}}, \\
    Q &= \frac{N_{\mathrm{test}}}{T_{\mathrm{test}}}
    \quad \text{flows per second}.
    \label{eq:inference_efficiency}
\end{align}

The automated Stage~2 profile measures inference after embeddings and context
indices have been constructed. It is therefore a \emph{model-evaluation}
measurement, not end-to-end PCAP-to-alert latency. The results chapter must
state this scope explicitly. If deployment latency is evaluated, extraction,
preprocessing, Stage~1 inference, context lookup, Stage~2 inference, and output
serialization are timed as separate components before reporting their sum.
Single-pass wall-clock values are treated as descriptive measurements; timing
noise should not be interpreted as a meaningful speed advantage when two
values differ only marginally.

\section{Implementation and Reproducibility}
\label{subsec:experimental_reproducibility}

The models are implemented in PyTorch \cite{Paszke2019PyTorch}, while
preprocessing and reported
classification metrics use scikit-learn
\cite{Pedregosa2011ScikitLearn}. Randomness is controlled by setting the Python,
NumPy, and PyTorch seeds. CUDA seeds are set for all devices,
\texttt{cudnn.deterministic} is enabled,
\texttt{cudnn.benchmark} is disabled, deterministic PyTorch algorithms are
requested, and \texttt{CUBLAS\_WORKSPACE\_CONFIG} is fixed. DataLoader workers
receive deterministic seeds derived from the loader generator.

At the start of each Colab run, the experiment log records the accelerator
name and memory, CPU model, available system RAM, Python version, PyTorch
version, CUDA runtime, cuDNN version, and scikit-learn version. These values are
captured from the active runtime rather than inferred from the Colab plan,
because a paid plan does not guarantee one fixed VM configuration
\cite{GoogleColabFAQ2026}.

Each run directory stores at least the following artifacts:
\begin{itemize}[leftmargin=*]
    \item the resolved YAML configuration and random seed;
    \item the split manifest, ordered flow identifiers, class counts, and
    preprocessing metadata;
    \item the selected feature lists and fitted preprocessing object;
    \item the best checkpoint, monitored validation score, and epoch history;
    \item validation and test probabilities, predicted labels, and selected
    validation threshold;
    \item the complete metric dictionary and confusion matrix;
    \item parameter counts, completed epochs, training time, inference
    profile, accelerator identity, GPU memory, and software versions; and
    \item for Stage~2, the exact Stage~1 checkpoint identifier, embedding-file
    checksums, context policy, relation, window size, and context-length
    distribution.
\end{itemize}

Embedding export uses non-sampled loaders so that the weighted training sampler
cannot create duplicate flows. Before Stage~2 training, exported metadata are
checked for unique flow identifiers, aligned array lengths, legal split names,
chronological order, and matching labels. These checks are part of the
experimental protocol rather than optional debugging because an embedding
misalignment can produce plausible but invalid performance.

The run manifest also records source-code revision, configuration checksum,
PCAP checksum, Suricata ruleset information, and custom extraction-plugin
revision where organizational policy permits. These practices follow broader
recommendations for transparent and reproducible machine-learning experiments
\cite{Pineau2021Reproducibility}.

\section{Result-Reporting Protocol}
\label{subsec:result_reporting_protocol}

The next chapter reports results in the same order as the experimental matrix.
For every predictive comparison, the main table contains the number of test
flows, macro-precision, macro-recall, macro-F1, class-1 precision, class-1
recall, class-1 F1, ROC-AUC, PR-AUC (AP), FPR, threshold, and the complete
confusion matrix. Efficiency tables add the accelerator type, trainable
parameters, completed epochs, training time, inference time, milliseconds per
flow, and flows per second.

The following reporting rules are applied:
\begin{enumerate}[leftmargin=*]
    \item Macro-F1 is the primary ranking metric; PR-AUC and class-1 metrics
    are used to interpret the ranking rather than to select a different winner
    after observing the test set.
    \item All decimal metrics are reported to four decimal places, thresholds
    to at least three decimal places, and confusion-matrix entries as integer
    counts.
    \item The exact test sample count accompanies every table. Results from
    different split manifests are not merged or compared by rank.
    \item The threshold shown for a test result is the threshold selected on
    validation predictions. Any test-oracle threshold is isolated in a
    diagnostic table and labelled non-deployable.
    \item The best value in a table is emphasized only among models evaluated
    on the same manifest and protocol. Statistical uncertainty and absolute
    effect size accompany claims of improvement.
    \item Negative or mixed results are retained. For example, a method that
    improves class-1 recall at the cost of FPR is described as an operational
    trade-off rather than an unconditional improvement.
\end{enumerate}

\section{Validity Boundaries}
\label{subsec:experimental_validity}

The protocol supports a controlled comparison on a large, chronologically
partitioned corporate capture, but it does not by itself establish universal
generalization. The labels are derived from Suricata alerts and may contain
false positives, false negatives, and rule-coverage bias. A single capture can
also encode organization-specific hosts, services, attack prevalence, and
collection conditions. The paired bootstrap quantifies uncertainty over the
observed test flows, not uncertainty over future organizations or rulesets.

Accordingly, conclusions are limited to the alert-derived binary detection task
and feature set described in Section~\ref{sec:dataset_preprocessing}. External
time-shifted or cross-network evaluation, when available, is reported separately
without refitting the preprocessor or threshold. This conservative scope is
consistent with warnings against treating closed-world NIDS benchmark accuracy
as direct evidence of deployment readiness
\cite{SommerPaxson2010,Arp2022DosDonts}.

\section{Chapter Summary}
\label{subsec:experimental_methodology_summary}

The experimental design evaluates the hierarchy in two controlled steps. The
first step tests packet-sequence, time-aware encoding, and flow-statistics
fusion within individual flows. The second uses one frozen Stage~1 embedding
set to test causal context selection and Stage~2 encoder alternatives. Model
selection and thresholding are confined to the validation partition, while the
chronologically latest flows form the held-out test set. Macro-F1 is the primary
metric, PR-AUC (implemented as average precision) measures rare-class ranking,
and FPR connects statistical performance to alert burden. Paired uncertainty
analysis and efficiency profiling complete the protocol, ensuring that the
results chapter can distinguish reliable architectural gains from threshold,
manifest, or computational differences.
